\chapter{After}

Recovery is two jobs. Rebuild trust in the keys. Look
after the humans the keys almost got killed.

\section{People}

Counseling is not a perk. Victims and families need it.

Ask, honestly, whether this person can still be a
key-holder. Their profile is now higher. Wanting to
``get back to normal'' is not an answer.

Tell the rest of the team what changed and what support
exists. Silence reads as ``we will burn the next signer
too.''

\section{Infrastructure}

\textbf{Do not reuse} seeds, keys, or devices that were
physically reachable, including the ones that ``look
fine.''

Build a new hierarchy. Move remaining assets only after
that hierarchy is real.

Re-check every signer, threshold, and allowlist. Incidents
are a good time to insert an attacker-controlled address.
Look.

\section{Governance}

Change the travel, remote-work, and high-value payment
rules that actually got used against you. Write the
lesson into the risk register and into onboarding so the
next hire is not told folklore.

\section{Watch}

Monitor the people who were hit, and the new wallets, for
follow-on surveillance or a second attempt.

Re-read the threat. The attacker who already succeeded
once has a cheaper second shot.

If you have physical-security or threat-intel partners,
this is when you keep them. Dropping the retainer after
the news cycle is how the sequel gets written.
